% ============================================================
%  THE ORTYX PROTOCOL
%  Part III — Decomposition
%  Chapter 14: The RSVP Audit
% ============================================================

\chapter{The RSVP Audit}
\label{ch:rsvp-audit}

\chapepigraph{A framework used to audit other frameworks
is not exempt from audit. If it were, it would not be
an audit framework. It would be a throne.}{Flyxion}

\noindent
Part~IV will use the \RSVPfull{} framework as a
reinterpretive substrate for the survivable components
of the Phoenix corpus. Before that deployment, RSVP
must pass through the same audit procedure that was
applied to the Phoenix corpus in Chapters~11--13.
This is not optional. A framework that audits others
but exempts itself from audit has not solved the
recursive closure problem; it has merely renamed it.

This chapter conducts the RSVP audit at the level of
detail appropriate for a reinterpretation framework
--- not the level required for a primary theoretical
framework, which would require a full separate monograph,
but the level required to establish that RSVP does
not replicate the failure geometries that disqualified
portions of the Phoenix corpus.

\section{RSVP in Brief}
\label{sec:rsvp-brief}

The Relativistic Scalar-Vector Plenum is a theoretical
framework developed by Flyxion that proposes a
field-theoretic account of the relationship between
physical structure, configurational accessibility,
and the emergence of observable phenomena. Its central
formal apparatus includes:

A scalar entropy field $\phi(x,t)$ representing the
accessibility of configurations at each point in
space-time. Regions of low $\phi$ correspond to
highly constrained, low-entropy configurations; regions
of high $\phi$ correspond to less constrained, higher-entropy
configurations.

A vector field $\mathbf{v}(x,t)$ representing the
flow of configurational possibility through the space,
which can be thought of as the ``current'' of admissible
trajectories.

An admissibility constraint operator $\mathcal{C}$
that specifies which field configurations are physically
realizable, and a dynamics that drives the fields
toward configurations consistent with the constraints.

A projection ladder connecting the field-theoretic
level to coarser-grained observational levels, with
each rung of the ladder corresponding to a different
level of description at which different constraints
become operative.

The framework is most directly applicable to physical
cosmology, where it has been developed as an alternative
to dark matter and dark energy explanations of
observational anomalies. Its interest for the present
project derives from its formal apparatus of projection
operators, accessibility entropy, and constraint
propagation, which are sufficiently general to serve
as a reinterpretive vocabulary for the Phoenix corpus's
survivable components.

\section{Applying the Audit Operators to RSVP}
\label{sec:rsvp-operators}

\subsection{Structural Consistency: \texorpdfstring{$\AuditS(\SysTriple_{\mathrm{RSVP}})$}{S(Theta\_RSVP)}}

RSVP's structural consistency score is highest for
claims derived from the field-theoretic apparatus.
The relationship between the scalar field $\phi$
and the vector field $\mathbf{v}$ is specified by
field equations that determine their joint evolution.
Claims about the dynamics of the coupled system
follow non-trivially from these equations.

The score drops for claims about how the RSVP framework
accounts for specific cosmological observations (galactic
rotation curves, lensing anomalies), where the derivation
requires additional assumptions about the field's
behaviour in specific regimes that are not uniquely
determined by the framework's axioms.

The score also drops for uses of RSVP as a general-purpose
reinterpretive vocabulary outside its primary cosmological
domain. When RSVP is applied to cognitive phenomena,
institutional dynamics, or epistemic processes --- as
it is in the reinterpretation work of Part~IV --- the
structural consistency of those applications depends
on the validity of analogical transport: whether the
same formal apparatus that is structurally grounded
in the cosmological context retains that grounding
when transposed to a different domain.

\begin{dangerzone}[Compression \(\neq\) Ontology]
RSVP's formal elegance and the breadth of its potential
applications create a specific risk: the projection
ladder and accessibility entropy concepts may be
applied to domains where they function as productive
metaphors rather than as structurally grounded
mechanisms. When RSVP is used to ``reinterpret'' Phoenix
concepts in Part~IV, the reinterpretation is productive
if it generates new research directions and clarifies
conceptual relationships; it risks becoming decorative
if the RSVP vocabulary is applied without establishing
that the quantities it names can be operationalized
in the new domain. The reinterpretation work of Part~IV
must be conducted with explicit acknowledgment of which
applications are structural and which are metaphorical.
\end{dangerzone}

\subsection{Falsifiability: \texorpdfstring{$\AuditF(\SysTriple_{\mathrm{RSVP}})$}{F(Theta\_RSVP)}}

RSVP's falsifiability score is highest for its
cosmological claims. Specific predictions about the
distribution of the scalar field in regions where
dark matter effects are observed, and about the
relationship between the vector field dynamics and
gravitational lensing patterns, are in principle
testable against astrophysical data. Whether the
current state of observational astronomy provides
the precision required to test these predictions
is an empirical question about the data, not about
the framework's structure.

The falsifiability score drops for applications outside
the cosmological domain. When RSVP concepts are applied
to cognitive or epistemic phenomena, specifying what
observations would disconfirm the application requires
specifying what RSVP predicts in those domains ---
which requires the operationalization discussed above.

\subsection{Projection Dependence: \texorpdfstring{$\AuditP(\SysTriple_{\mathrm{RSVP}})$}{P(Theta\_RSVP)}}

RSVP's most significant audit vulnerability is precisely
the projection ladder that makes it useful as a
reinterpretive framework. The ladder provides a formal
apparatus for moving between levels of description,
but the constraint-faithfulness of each step in the
ladder --- whether the constraints operative at the
fine-grained level are preserved by the projection
to the coarser level --- requires establishment for
each specific application.

In the cosmological context, constraint-faithfulness
at some rungs of the ladder is established by standard
renormalization group arguments; at other rungs, it
is an open theoretical question. In cognitive and
epistemic applications, the constraint-faithfulness
of the projection is less well-established, and
treating it as given would constitute the same
projection dependence vulnerability identified in
the Phoenix corpus.

\begin{auditprop}
\label{ap:rsvp-projection}
The RSVP framework's projection ladder is its most
significant audit vulnerability in its current state.
The ladder provides structural legitimacy for
applications that remain close to its cosmological
grounding; it provides metaphorical resources for
applications in other domains; but the constraint-faithfulness
of applications across domain boundaries is not
established by the framework's own apparatus. Using
RSVP as a reinterpretive framework in Part~IV therefore
requires explicit acknowledgment at each step of which
applications are within the domain of established
constraint-faithfulness and which are metaphorical
extensions.
\end{auditprop}

\subsection{Recursive Closure: \texorpdfstring{$\AuditR(\SysTriple_{\mathrm{RSVP}})$}{R(Theta\_RSVP)}}

RSVP's recursive closure risk is moderate and structurally
specific. The framework's primary closure risk arises
from the generality of the accessibility entropy concept.
Like the Phoenix coherence concept, accessibility entropy
is defined precisely in its original domain (the logarithm
of the size of the admissible trajectory space) and
can be applied analogically in other domains. If the
analogical applications are not constrained by
operationalization requirements, the accessibility
entropy concept could undergo the same cascade of
broadening that the coherence concept underwent in
the Phoenix corpus.

The framework has one structural feature that partially
mitigates this risk: the admissibility constraint
operator $\mathcal{C}$ provides an independent handle
on what counts as an admissible trajectory, separate
from the accessibility entropy. This means there is
a second concept that must be operationalized alongside
the entropy, providing a check against unconstrained
broadening of the entropy concept alone.

Whether this structural feature is sufficient to prevent
$\FailGeo{3}$ from developing in extended applications
of RSVP is an open question. The risk is real and the
mitigation is partial.

\section{The RSVP Audit Result}
\label{sec:rsvp-result}

Applying the four operators, the RSVP audit can be
summarized:

\begin{itemize}
  \item $\AuditS$: High within the cosmological domain;
  moderate to low in cross-domain applications.
  \item $\AuditF$: High for cosmological predictions;
  conditional on operationalization for other applications.
  \item $\AuditP$: Elevated; the projection ladder is
  a vulnerability as well as a resource.
  \item $\AuditR$: Moderate; accessibility entropy faces
  a broadening risk analogous to the coherence cascade.
\end{itemize}

The RSVP audit result is better than the Phoenix audit
result, which is why RSVP is deployed as the reinterpretive
framework in Part~IV rather than Phoenix. But it is
not a clean result: RSVP has real vulnerabilities that
Part~IV's reinterpretation work must navigate explicitly.

\begin{auditprop}
\label{ap:rsvp-qualified}
The RSVP framework passes the \Ortyx{} audit at a
level sufficient to function as a reinterpretive
substrate, with the following explicit qualifications:
(1)~RSVP applications in Part~IV are structural
where the projection is to the cosmological domain
or where operationalization is established; they
are metaphorical elsewhere, and must be labeled as
such. (2)~The accessibility entropy concept must
be applied with operationalization constraints that
prevent the cascade of broadening identified in the
Phoenix coherence concept. (3)~The projection
ladder's constraint-faithfulness must be acknowledged
as unestablished for cognitive applications, treating
it as a productive research hypothesis rather than
an established result.
\end{auditprop}

\section{Why RSVP Wins}
\label{sec:why-rsvp-wins}

The question of why RSVP rather than some other
framework serves as the reinterpretive substrate
deserves an explicit answer. RSVP is not deployed
here because it is ``true'' in some absolute sense
that the present analysis cannot establish. It is
deployed because it exhibits lower projection defect
for the specific survivable components of the Phoenix
corpus than the available alternatives.

The projection defect comparison works as follows.
The Phoenix corpus's survivable components include
the jitter metric, the conformity relation, the wave
superposition memory efficiency, and the structural
parallel to predictive processing. Each of these
can be translated into RSVP-compatible language with
specifiable constraint-faithfulness:

The jitter metric translates to accessibility entropy
in the signal's state space: low jitter corresponds
to low-entropy (highly constrained) trajectory regions.
The translation is metaphorical but generates a
specific research direction (testing whether jitter
correlates with trajectory space constraint metrics).

The conformity relation translates to constraint
propagation: harmonics that conform to the fundamental's
envelope are evolving within a shared admissible region.
The translation is productive and suggests connections
to the RSVP constraint dynamics.

The wave superposition memory efficiency translates
naturally to field superposition in the RSVP plenum.
This is the tightest translation, as both formalisms
involve wave interference in a field.

The predictive processing parallel translates to
constraint-governed trajectory selection: the system
selects trajectories consistent with its generative
model of the world, and high-salience events are those
that require trajectory revision.

These translations are not perfect. But they are
better-grounded than the available alternatives,
and ``better-grounded'' is all that the \Ortyx{}
Protocol requires for reinterpretation substrate
selection.

\medskip

Chapter~15 synthesizes the Part~III decomposition
into a summary audit score for the Phoenix corpus
and prepares the classification of survivable,
unstable, and metaphorically useful components
that will structure the Interlude.
